\section{Rome and the Manner of Death}

No New Testament text says that Peter went to Rome or describes his death.
The case for both is cumulative, and a cumulative case is only worth as much
as it survives of the case against it. This section therefore states the case
against first, at full strength, and then asks what the surviving first- and
second-century evidence actually carries.

\subsection{What Scripture gives}

Within Scripture there are two pointers, and neither is a statement. The
first is the risen Christ's prophecy that Peter, grown old, would ``stretch
forth thy hands'' while ``another shall gird thee,'' written down with the
evangelist's comment that it signified ``by what death he should glorify
God'' (Jn 21:18--19). This establishes that by the end of the first century
Peter's death was known to have been a martyrdom, and that the Johannine
circle expected its readers to know it too. It names no place, no manner, and
no date. The second is the greeting of 1~Peter from ``Babylon'' with Mark
(1~Pt 5:13), read as Rome by the second-century tradition. Neither is proof
by itself; both were written when the outcome was already memory.

\subsection{The case against, at its strongest}

Four silences make a serious argument, and they are usually stated too
weakly by those who reject the tradition and too dismissively by those who
hold it.

\textbf{The silence of Romans.} Paul wrote to the church at Rome about
A.D.~57, before he had ever visited it. The letter's last chapter greets some
two dozen Roman Christians by name---Prisca and Aquila, Epaenetus, Mary,
Andronicus and Junia, Ampliatus, Urbanus, Apelles, Herodion, Tryphaena and
Tryphosa, Persis, Rufus and his mother, and a further list of names in
households---and never mentions Peter. Paul also states, in the same letter,
the principle governing his own mission: ``I have so preached this gospel,
not where Christ was named, lest I should build upon another man's
foundation'' (Rom 15:20). On his own stated rule, a Roman church already
founded and led by Peter is a church Paul would not have written to in this
way, and a Peter present in Rome in 57 is a man Paul could not have omitted
from that list of greetings.

\textbf{The silence of the captivity letters.} The letters traditionally
assigned to a Roman imprisonment name Paul's companions there repeatedly.
Colossians lists Aristarchus, Mark, Jesus called Justus, Epaphras, Luke, and
Demas; Philemon lists Epaphras, Mark, Aristarchus, Demas, and Luke; 2 Timothy
says ``only Luke is with me'' and asks for Mark. Peter appears in none of
them.

\textbf{The silence of Acts.} Luke's second volume is structured to bring the
gospel from Jerusalem to Rome. It ends with Paul in Rome, preaching for two
years in his own hired lodging. If Peter had died there in the same
persecution, or had ministered there at all, the book that is otherwise so
interested in Peter says nothing whatever about it.

\textbf{The silence of the first century as a whole.} No surviving text from
the first century---not a Gospel, not a letter, not Acts, not Revelation,
not Josephus or any Roman writer---places Peter in the city of Rome.

Taken together these are not trivial. They are decisive against one version
of the tradition: the version, developed by the third century and transmitted
by Jerome, in which Peter arrived in Rome in the second year of Claudius and
held the see for twenty-five years. That version is refuted by Galatians and
Acts alike, which have Peter in Jerusalem and Antioch through the late
forties, and it is refuted a second time by Romans 16, which has no Peter in
Rome in 57. This study rejects it, and rejects it on exactly this evidence.

\subsection{Why the silences do not reach the claim actually made}

The claim the early sources make is narrower than the claim the silences
refute. It is that Peter came to Rome late, taught there, and died there
under Nero---within a window that opens after Paul wrote Romans and closes
with Nero's death in 68. Against that claim the silences are much weaker.

Romans and the captivity letters bear on Peter's whereabouts in the years
they were written, which is precisely the period the tradition does not
require him to be in Rome. The argument from Romans 15:20 cuts the other way
as well: if Paul thought Peter had founded the Roman church, he says he would
not have written; he did write; therefore, on his own principle, Peter had
not founded it. That is a conclusion about foundation, not about death, and
it is one this study accepts. When Irenaeus later calls the Roman church
``founded and organized\ldots\ by the two most glorious apostles, Peter and
Paul,'' he is making a claim about the church's apostolic ordering and cult,
not a claim about who first preached there; Christians existed at Rome before
either apostle arrived, as Paul's own letter proves.

The silence of Acts is the weakest of the four, because it can be shown to
be selective rather than ignorant. Acts omits the Antioch incident, which
Paul reports in his own hand and which no one doubts; Jerome, arguing with
Augustine about that very passage, states the principle without noticing that
he is stating it---``although the Acts of the Apostles do not mention this,
but we must believe Paul's statement.'' Acts also omits the death of Paul,
its own protagonist, though it takes him to Rome and leaves him there under
guard. A book that declines to narrate the death of the man it has followed
for sixteen chapters is not a book from whose silence one may infer that
another man did not die. That is not a modern rationalization: the Muratorian
Fragment, about 170, gives exactly this account of Luke's method, saying that
he compiled the events that took place in his presence, ``as he plainly shows
by omitting the martyrdom of Peter as well as the departure of Paul from the
city.'' Whatever one thinks of the fragment's explanation, its author knew
both that Peter had been martyred and that Acts does not say so.

\subsection{The evidence for Rome, weighed by date}

\textbf{1 Clement, from Rome, about 95--96.} The earliest external witness is
the letter of the Roman church to Corinth known as 1~Clement, and it repays
being read in its context rather than quoted in a clause. Its fourth chapter
runs through a series of Old Testament examples of what envy and jealousy
have destroyed---Cain, Jacob, Joseph, Moses, Aaron and Miriam, Dathan and
Abiram, David. Its fifth chapter then turns to the present: ``But not to
dwell upon ancient examples, let us come to the most recent spiritual heroes.
Let us take the noble examples furnished in our own generation.'' Peter comes
first: ``Peter, through unrighteous envy, endured not one or two, but
numerous labours; and when he had at length suffered martyrdom, departed to
the place of glory due to him.'' Paul follows, with a longer notice ending
``having taught righteousness to the whole world, and come to the extreme
limit of the west, and suffered martyrdom under the prefects.'' Then chapter
six adds ``a great multitude of the elect, who, having through envy endured
many indignities and tortures, furnished us with a most excellent example''---and among them ``those women, the Danaids and Dircae, being persecuted, after
they had suffered terrible and unspeakable torments, finished the course of
their faith.''

Three features of that passage matter. The letter names no city and no manner
of death for Peter, and this study does not pretend otherwise. But it is the
church \emph{of Rome} writing, it calls these martyrdoms the examples of
``our own generation,'' and it sets Peter and Paul at the head of a local
mass martyrdom in which women are described under the names of figures from
Greek myth, ``the Danaids and Dircae.'' That those names most probably refer
to executions staged as mythological tableaux is this study's reading of the
clause rather than something the letter says; what the letter does say is
that these women suffered ``terrible and unspeakable torments'' as public
spectacle. Tacitus describes precisely such staging in the aftermath of the
fire of 64: an ``immense multitude'' convicted; ``mockery of every
sort was added to their deaths. Covered with the skins of beasts, they were
torn by dogs and perished, or were nailed to crosses, or were doomed to the
flames and burnt, to serve as a nightly illumination'' in Nero's gardens
(\work{Annals} 15.44). The Roman letter's ``great multitude'' and the Roman
historian's ``immense multitude'' are describing the same kind of event in
the same city within a generation of each other, and one of them puts Peter
at the head of the list.

\textbf{Ignatius, to Rome, about a decade later.} Writing ahead to the church
of Rome on his way to execution, Ignatius uses the two apostles as the
measure of an authority he disclaims: ``I do not, as Peter and Paul, issue
commandments unto you. They were apostles; I am but a condemned man''
(\work{Romans} 4). The sentence is worth more than its length. It is
addressed to the Roman church specifically, it takes for granted that Peter
and Paul had stood in a position to command that church, and it comes from
Antioch---a second city, independently, treating the Roman association as
common knowledge.

\textbf{Dionysius of Corinth, about 170.} Writing to Rome, he says that Peter
and Paul ``planted and likewise taught us in our Corinth. And they taught
together in like manner in Italy, and suffered martyrdom at the same time''
(in Eusebius, \work{Church History} 2.25.8). This is a bishop of one apostolic
church writing to another about a shared memory; his phrase about the timing
is loose enough to mean the same period rather than the same day, and it is
so treated in the tradition audit.

\textbf{Irenaeus, about 180.} At Lyons, with Asian roots, he speaks of ``the
very great, the very ancient, and universally known Church founded and
organized at Rome by the two most glorious apostles, Peter and Paul,'' and
begins the Roman succession from their committing the episcopate to Linus
(\work{Against Heresies} 3.3.2--3); in the previous book he dates Matthew's
Gospel ``while Peter and Paul were preaching at Rome'' (3.1.1).

\textbf{Gaius, about 200, pointing at monuments.} The most concrete witness
is the least rhetorical. Eusebius reports that Gaius, ``a member of the
Church, who arose under Zephyrinus, bishop of Rome''---which dates the
dialogue between 199 and 217---in a published disputation with Proclus, the
leader of the Phrygian heresy, wrote: ``But I can show the trophies of the
apostles. For if you will go to the Vatican or to the Ostian way, you will
find the trophies of those who laid the foundations of this church''
(\work{Church History} 2.25.6--7). This is not a report of a tradition; it is
a challenge to an opponent in the same city to walk out and look. Whatever
Gaius meant by \textit{tropaion}---a monument, a memorial, a
victory-marker---he expected Proclus to be able to find two of them, at two
identified places, and he expected Proclus to accept that they marked the
graves of Peter and Paul. The polemical setting is the point: a disputant
does not stake an argument on a monument his opponent can check unless he is
confident of what will be found there.

\textbf{The absence of a rival.} Finally, and this is a negative that
functions as evidence: no other church in antiquity ever claimed Peter's
death or his tomb. Jerusalem, Antioch, and Corinth each had a serious Petrine
association and each had every incentive; none of them made the claim. In a
world of competing apostolic prestige, that unanimity is difficult to
manufacture.

\subsection{Manner and date}

Crucifixion is early, credible, and not contemporary. The first explicit
statement is Tertullian's, about 200: Rome is the church ``where Peter
endures a passion like his Lord's! Where Paul wins his crown in a death like
John's'' (\work{Prescription} 36); and again, of Nero's persecution, ``Then
is Peter girt by another, when he is made fast to the cross''
(\work{Scorpiace} 15)---a clause that is itself an echo of John 21:18, which
is a caution against treating it as independent information. It is compatible
with Tacitus's crosses and with the Johannine prophecy, and no ancient source
contradicts it; it is not first-century testimony.

The head-downward detail appears a generation later still. Eusebius reports
from Origen, about 230, that Peter ``was crucified head-downwards; for he had
requested that he might suffer in this way'' (\work{Church History} 3.1.2). It
had already been narrated, with a quite different motive, in the apocryphal
\work{Acts of Peter} treated in the next section.

The year is unrecoverable. The persecution opens in 64; the chronicle
tradition used by Jerome put the death in Nero's fourteenth year, 67 or 68.
``Between 64 and 68'' is all the evidence yields, and even that depends on
accepting that the death belongs to the Neronian persecution rather than
merely to Nero's reign.

\begin{studybox}{Project synthesis --- the Roman death}
\textbf{Judgment.} That Peter died as a martyr at Rome in the 60s is
historically well attested and morally certain in Catholic tradition. This
judgment is not the same as strict documentary certainty: no surviving
first-century text explicitly states the city, date, and manner of his death.
That he ever governed the Roman church as its bishop for a term of years is
not established by any evidence this study could reach, and the
twenty-five-year episcopate is positively excluded.

\textbf{Reasons.} (1) A letter from the Roman church itself, within roughly
thirty years of the event, names Peter first among the martyrs of ``our own
generation'' and sets him beside a local mass martyrdom whose details match
what Tacitus reports of Nero's spectacles. (2) An independent letter from
Antioch, a decade later, addressed to Rome, assumes that Peter had authority
there. (3) Corinth in 170, Lyons in 180, and Rome itself in 200 agree, and
the last of them points at visitable monuments in a hostile debate. (4) No
competing city ever claimed the death or the grave. (5) By contrast, every
element of the long-episcopate version collides with Galatians, Acts, and
Romans 16, and this study discards it.

\textbf{Strongest counterargument.} That the whole chain is one tradition
propagating, not several witnesses converging. 1~Clement names no city, so
its evidential value depends entirely on an inference from where it was
written; Ignatius names no city either and could be repeating what he had
heard; Dionysius, Irenaeus, and Gaius all write after the Roman church had
acquired an institutional interest in the claim, and all three could be
downstream of a single second-century Roman conviction whose origin is
invisible to us. On this account the tradition is early but not independent,
and early unanimity about a claim that only one interested party was in a
position to make is exactly what one would expect whether or not it were
true.

\textbf{What would change this judgment.} A first-century or early
second-century text placing Peter's death anywhere else; or a demonstration
that 1~Clement is not of Roman provenance or is substantially later than the
conventional date, which would remove the one witness that stands within
living memory. Conversely, a securely dated first-century witness naming Rome
explicitly would supply the strict documentary certainty that the present
cumulative case does not. None of these is known to this study.
\end{studybox}
